\documentclass[11pt]{article}
\usepackage[margin=1in]{geometry}
\usepackage{amsmath,amssymb,amsthm,mathtools}
\usepackage[T1]{fontenc}\usepackage{lmodern}\usepackage{microtype}
\usepackage{fancyhdr}
\usepackage{hyperref}
\setlength{\headheight}{14pt}
\pagestyle{fancy}
\fancyhf{}
\fancyhead[L]{\small DRAFT --- UNVERIFIED}
\fancyhead[R]{\small 6 October 2026}
\fancyfoot[C]{\thepage}
\hypersetup{colorlinks=true,linkcolor=blue,citecolor=blue,urlcolor=blue,
  pdftitle={Kirchberg algebras as full amalgams of stably finite algebras -- Unverified draft},
  pdfauthor={Caleb Barnett},
  pdfsubject={Proposed proof; public draft 1; not independently verified}}
\newtheorem{theorem}{Theorem}
\newtheorem{lemma}[theorem]{Lemma}
\newtheorem{proposition}[theorem]{Proposition}
\theoremstyle{remark}\newtheorem{remark}[theorem]{Remark}
\newcommand{\Z}{\mathbb Z}\newcommand{\C}{\mathbb C}
\newcommand{\T}{\mathbb T}\newcommand{\KK}{\operatorname{KK}}
\newcommand{\coker}{\operatorname{coker}}\newcommand{\rank}{\operatorname{rank}}
\newcommand{\Her}{\operatorname{Her}}\newcommand{\id}{\mathrm{id}}
\title{Kirchberg algebras as full amalgams\\
of stably finite algebras}
\author{Caleb Barnett}
\date{6 October 2026 --- Public draft 1}
\begin{document}\maketitle\thispagestyle{fancy}
\begin{center}
\begin{minipage}{0.92\textwidth}
\small\textbf{Draft --- unverified.} This manuscript presents a proposed
proof. It has not been independently verified by a human mathematician
and may contain errors. It is shared for scrutiny and comment.
\end{minipage}
\end{center}
\begin{abstract}
We present a proposed proof that every separable nuclear unital simple
purely infinite C*-algebra $B$ admits a decomposition
$B\cong A_0*_D A_1$, where $D,A_0,A_1$ are separable nuclear simple
unital stably finite C*-algebras and the common inclusions are unital
and injective. The product carries the full universal norm. No UCT
assumption is made on $B$. The construction combines finite-dimensional
ordinary-graph amalgams with a nuclear residually finite-dimensional
coefficient algebra KK-equivalent to $B$. Three positive embedding
channels preserve the reduced coefficient KK-class, cancel its scalar
summand, and enforce simplicity. A projection in the common algebra
prescribes the unit, and unit-preserving KK classification supplies the
final identification. If correct, the argument answers the question in
Suzuki's Remark~1.5 for all unital Kirchberg algebras.
\end{abstract}
\section{Statement and conventions}
\begin{theorem}[Proposed main theorem]\label{thm:all}
Every unital Kirchberg algebra $B$ has a unital isomorphism
\[
                     B\cong A_0*_D A_1
\]
where $D,A_0,A_1$ are separable nuclear simple unital stably finite
C*-algebras and both inclusions $D\to A_j$ are unital and injective.
The product is the full unital amalgamated free product.
\end{theorem}
A Kirchberg algebra means separable, nuclear, simple and purely infinite;
no UCT assumption is implicit. All tensor products are minimal except
where the maximal universal property is explicitly used. Matrix
multiplicity matrices act on column vectors, with target blocks as rows.
A minimal matrix projection is the $K_0$ generator. $SX=C_0((0,1))\otimes X$.
We write Kasparov composition from left to right: $x\otimes_Y y$ first
applies $x\in\KK(X,Y)$ and then $y\in\KK(Y,Z)$.

The question is posed in \cite[Remark 1.5, p.~2]{Suzuki}; the constituent
requirements are those of its Corollary 1.4. Nuclearity of the three
constituents is an additional conclusion here. No expectation onto the
limit common algebra is required or asserted. The proof uses UCT only
for explicitly identified finite graph algebras, never for $B$ or the
coefficient algebra.

\section{An ordinary graph as a full finite-dimensional amalgam}
Let $A$ be an $r\times r$ nonnegative integer matrix with no zero column,
and let $k,\ell\in\Z^r$ have strictly positive entries. Define
\begin{align*}
 D(A;k,\ell)&=\bigoplus_{v=1}^rM_{k_v}\ \oplus\ \bigoplus_{v=1}^rM_{\ell_v},\\
 F_0(A;k,\ell)&=\bigoplus_vM_{k_v+\ell_v},\\
 F_1(A;k,\ell)&=\bigoplus_vM_{k_v+\sum_w A_{vw}\ell_w}.
\end{align*}
The common inclusions have multiplicity matrices $[I\ I]$ and $[I\ A]$.
Write $P(A;k,\ell)=F_0*_{D}F_1$.

Let $E_A$ have $A_{vw}$ edges from $w$ to $v$; its usual source-row
adjacency matrix is $A^t$. Write its vertex projections $p_v$ and edge
partial isometries $s_{wvt}$, where $1\le t\le A_{vw}$. Thus
\[
 s_{wvt}^*s_{wvt}=p_v,\qquad
 \sum_{v,t}s_{wvt}s_{wvt}^*=p_w.
\]
Distinct edges have orthogonal range projections. With $m=k+\ell$ and
$M=\sum_vm_v$, put $q=\operatorname{diag}(p_v\text{ repeated }m_v\text{ times})$
in $M_M(C^*(E_A))$.

\begin{lemma}\label{lem:graph}
There is a unital isomorphism
\[
                  P(A;k,\ell)\cong qM_M(C^*(E_A))q.
\]
The projection $q$ is full. The minimal projection of the $v$th
$F_0$ summand corresponds to $[p_v]$ in $K_0$, and the product unit
corresponds to the vector $k+\ell$.
\end{lemma}
\begin{proof}
Index the $v$th $F_0$ matrix units $E^v$ by $P_{v,a}$ and $Q_{v,b}$,
with $a\le k_v$ and $b\le\ell_v$. Index the $v$th $F_1$ units $F^v$
by $P_{v,a}$ and $Q_{w,t,b}$, with $t\le A_{vw}$, $b\le\ell_w$.
The common relations are
\[
 E^v_{P_{v,a},P_{v,b}}=F^v_{P_{v,a},P_{v,b}},\qquad
 E^w_{Q_{w,b},Q_{w,c}}=\sum_{v,t}F^v_{Q_{w,t,b},Q_{w,t,c}}.
\]
In the proposed corner send $E^v_{ij}$ to $e_{ij}\otimes p_v$.
Writing $a_v=P_{v,1}$, define
\[
 V^v_{P_{v,a}}=e_{P_{v,a},a_v}\otimes p_v,\qquad
 V^v_{Q_{w,t,b}}=e_{Q_{w,b},a_v}\otimes s_{wvt}.
\]
Within each $v$ these have common initial projection
$e_{a_v,a_v}\otimes p_v$ and orthogonal final projections. Final
projections for different $v$ are also orthogonal. Their total sum is
$q$ by the graph relations. Hence $V^v_\alpha(V^v_\beta)^*$ defines
a unital representation of $F_1$, agreeing with that of $F_0$ on $D$.
The full universal property gives a homomorphism $\Psi$ to the corner.

Conversely put $p_v^P=E^v_{a_v,a_v}$, $p=\sum_vp_v^P$, and
\[
 t_{wvt}=E^w_{P_{w,1},Q_{w,1}}F^v_{Q_{w,t,1},P_{v,1}}\in pPp.
\]
Matrix multiplication gives the graph relations, with initial projection
$p_v^P$, orthogonal ranges for distinct edges, and range sum $p_w^P$
over edges starting at $w$. The graph universal property
\cite[Theorem 1.2]{KPR} gives $\theta:C^*(E_A)\to pPp$.
For a matrix entry $x_{ij}\in p_vC^*(E_A)p_w$ set
\[
 \Theta([x_{ij}])=\sum_{i,j}E^v_{i,a_v}\theta(x_{ij})E^w_{a_w,j}.
\]
The orthogonal matrix units make this a *-homomorphism from the corner.
Both composites fix the $E$ units. They also fix the graph edges in the
chosen vertex corner and the mixed units
\[
 F^v_{Q_{w,t,b},P_{v,a}}
       =E^w_{Q_{w,b},P_{w,1}}t_{wvt}E^v_{P_{v,1},P_{v,a}}.
\]
These mixed units, their adjoints and products recover all $F$ units;
the graph edges and the $E$ units generate the inflated graph corner.
Thus both composites are identities. Each vertex occurs in $q$, so $q$
is full. The generator and unit statements follow from these explicit maps.
\end{proof}

We henceforth set $A=I+B$, where every entry of $B$ is a positive integer.
The finite graph is strongly connected, has no sinks, and has at least
two loops at every vertex. It is cofinal and every cycle has an exit.
Its algebra is simple and purely infinite by \cite[Corollary 3.11]{KPR},
and nuclear by the ordinary graph nuclearity statement in
\cite[Remark 3.10]{AG}. The full corner in Lemma \ref{lem:graph} is
therefore a unital Kirchberg algebra. No separated-graph nonnuclear
full product is substituted for this ordinary graph algebra.

Finite-dimensional inclusions have unital conditional expectations.
The sequence from \cite[Theorem 4.1]{FG} gives the integer map
\[
 \phi(x,y)=(x+y,-x-Ay),\qquad
 K_0(P)=\coker B,\quad K_1(P)=\ker B,\quad [1_P]=[k+\ell].
\]
We identify $y\in\ker B$ with $(-y,y)\in\ker\phi$. The $K_0$ generators
are the minimal projections of $F_0$. These conventions will fix all
bonding maps, rather than merely the abstract groups.

\section{Lifting arbitrary graded maps by positive diagrams}
Let source and target matrices be $B=A-I$ and $B'=A'-I$, of sizes
$r$ and $r'$. Suppose integer matrices $S,H$ satisfy $SB=B'H$.
For any $r'\times r$ integer matrix $Z$, define
\begin{equation}\label{eq:T}
 T=\begin{pmatrix}S-Z&S-H-ZA\\Z&H+ZA\end{pmatrix},
 \qquad R=S+B'Z.
\end{equation}
Direct multiplication gives
\begin{equation}\label{eq:squares}
 [I\ I]T=S[I\ I],\qquad [I\ A']T=R[I\ A].
\end{equation}
When $T,S,R$ are strictly positive, set
$\binom{k'}{\ell'}=T\binom{k}{\ell}$. These identities give actual
unital finite-dimensional embeddings of multiplicities $T,S,R$.
Fix the standard common-algebra embedding. The two representations of
$D$ in each target factor have equal irreducible multiplicities; match
equal copies with a permutation unitary and conjugate the whole factor
embedding. The squares then commute literally. The induced full-product
homomorphism is injective because its source is simple and it is unital.

The maps on $K$-theory are $[S]$ on the cokernel and $H$ on the kernel.
Indeed the latter follows directly from
$T(-y,y)=(-Hy,Hy)$ when $By=0$, and the former from the first-factor
minimal projections. This also follows from naturality of the cone sequence.

\begin{lemma}[Simultaneous positive lifting]\label{lem:positive}
Let $B,B_0'$ be strictly positive square integer matrices, with target
size at least two. For any finite list of prescribed pairs
\[
 f_{0,c}:\coker B\to\coker B_0',\qquad
 f_{1,c}:\ker B\to\ker B_0',
\]
one can replace the target presentation by $B'=UB_0'$, with $U$
nonnegative unimodular, and find strictly positive matrices $T_c,S_c,R_c$
as in \eqref{eq:T}, realizing each pair under the induced target
identifications. All channels use the same $B'$.
\end{lemma}
\begin{proof}
First lift $f_{0,c}$ to an integer matrix $S_{0,c}$ on the free presentations.
The columns of $S_{0,c}B$ lie in $\operatorname{im}B_0'$, so choose
$H_{0,c}$ with $S_{0,c}B=B_0'H_{0,c}$. The difference
$f_{1,c}-H_{0,c}|_{\ker B}$ lands in $\ker B_0'$. Since
$\Z^r/\ker B\cong\operatorname{im}B$ is free, $\ker B$ is a direct
summand of $\Z^r$. Extend the difference to a map into $\ker B_0'$
and add it to $H_{0,c}$. This realizes both prescribed maps.

Let $J$ denote an all-ones matrix of the appropriate rectangular size.
For sufficiently large positive integers $t_c$, the matrices
\[
 \overline S_c=S_{0,c}+B_0'(t_cJ),\qquad
 \overline H_c=H_{0,c}+(t_cJ)B
\]
are strictly positive. These additions preserve the chain equation and
both induced group maps. Set $Z_c=J$ and
$W_c=\overline H_c+J A$. Choose a common nonnegative unimodular $U$
whose every row sum exceeds every entry of every $W_c$ and exceeds one.
Such matrices exist in size $r'\ge2$: add a large multiple of row two
to row one of the identity, and then a large multiple of the new row one
to every other row. Each operation has determinant one.

Put $B'=UB_0'$, $S_c=U\overline S_c$, and $H_c=\overline H_c$.
Every entry of $S_c$ is at least the corresponding row sum of $U$.
Thus $S_c-J$ and $S_c-W_c$ are strictly positive. The other blocks
$J,W_c$ and $R_c=S_c+B'J$ are strictly positive as well.
Left multiplication by $U$ changes neither the kernel nor the cokernel
up to the stated identification and does not alter either prescribed map.
\end{proof}

For multiple channels, the dimensions are defined with
$T_{\mathrm{tot}}=\sum_cT_c$, not with any single $T_c$.
This gives a useful nonunital version of the construction. In each target
common summand put the different channel blocks in disjoint coordinate
subspaces. Their units are orthogonal projections $q_c\in D'$ with
$\sum_cq_c=1$. The channel factor maps have units $i_j'(q_c)$.
Their restrictions to $D$ match the chosen channel embedding, by the same
matrix representation equivalence as above. Hence the universal maps are
\[
                    \gamma_c:P\longrightarrow q_cP'q_c.
\]
They are nonzero and injective, and their ranges are orthogonal. Each
$q_c$ is full in $D'$ because every channel multiplicity is positive.
Their sum is a unital product map. The cone maps and boundary maps are
natural also for these nonunital homomorphisms, so the individual channel
maps still induce $S_c,H_c$ as above. No compatibility of conditional
expectations is required for this naturality: the cone map is the
pointwise map on its defining continuous functions.


\section{A nuclear residually finite-dimensional KK model}
An algebra is residually finite-dimensional (RFD) if finite-dimensional
*-representations separate its points. A separable RFD algebra has a
countable separating family: select representations detecting a countable
dense subset, with arbitrarily good norm bounds, and take their union.
RFD passes to subalgebras and to minimal unitization. No UCT hypothesis is assumed or used for the nuclear RFD
coefficient algebra.

\begin{lemma}\label{lem:rfd}
For every separable nuclear algebra $B$, there is a separable nuclear
RFD algebra $F$ and an invertible $\xi\in\KK(F,B)$.
\end{lemma}
\begin{proof}
The construction in \cite[Lemma 2.4 and proof of Theorem 1.2]{Dadarlat}
gives a separable nuclear RFD extension
\[
       0\longrightarrow I=\bigoplus_nM_{r(n)}
       \xrightarrow{\theta} H\longrightarrow SB\longrightarrow0.
\]
For clarity, apply its quasidiagonal pullback construction to
$\widetilde{SB}$, obtaining an algebra inside $\prod_nM_{r(n)}$,
and take the inverse image of $SB$. The quotient is nuclear, so the
extension is semisplit; the ideal and quotient are nuclear, so $H$ is
nuclear. It is RFD as a subalgebra of the matrix product.

Let
\[
 F=C_\theta=\{(f,x)\in C_0([0,1),H)\oplus I:f(0)=\theta(x)\}.
\]
The extension $0\to SH\to F\to I\to0$ proves nuclearity.
The two ambient summands are RFD, so $F$ is RFD. It is separable.
The mapping-cone triangle of the semisplit extension above identifies
$C_\theta$ up to KK-equivalence with $S(SB)$. Bott periodicity supplies
$\xi$. The hypothesis $K_*(B)=0$ used later in Dadarlat's proof is not
used in this construction or mapping-cone equivalence.
\end{proof}

Fix $F$ and $\xi$ for the target $B$. Put $E=\widetilde F$, with
inclusion $\iota:F\to E$, scalar quotient $\varepsilon:E\to\C$,
and scalar section $\eta:\C\to E$. Use the external unitization if
necessary, so that $0\to F\to E\to\C\to0$ is always split.
Split exactness gives a class $r\in\KK(E,F)$ such that
\begin{equation}\label{eq:split}
 [\iota]\otimes_Er=1_F,\qquad
 r\otimes_F[\iota]=e:=1_E-[\varepsilon]\otimes_\C[\eta].
\end{equation}
In particular $e^2=e$. Choose unital finite-dimensional representations
\[
                 \sigma_h:E\longrightarrow M_{q_h}(\C)
\]
whose every tail separates points. For example, repeat each member of
a countable separating family infinitely often. Their direct sum is
faithful, so $\{\id_X\otimes\sigma_h\}$ separates $X\otimes E$ for
every $X$: represent $X$ faithfully and use the block decomposition of
the faithful spatial representation. No homotopy condition on $\sigma_h$
is needed.

\section{The three actual coefficient channels}\label{sec:channels}
We use only graph presentations with $K_0=\Z$ and $K_1=\Z$.
Start with
\[
 B_1=\begin{pmatrix}1&1\\1&1\end{pmatrix},\qquad A_1=I+B_1.
\]
The cokernel identification sends $[e_1]$ to $1$ and $[e_2]$ to $-1$;
the kernel has generator $(1,-1)$. At later steps take a positive
presentation $B_{h+1}=U_hB_1$ provided by Lemma \ref{lem:positive},
with the induced cokernel identification and the same kernel convention.
The source is the previously chosen $B_h$; there is no requirement that
$U_h$ relate consecutive presentations directly.

Write $D_h,F_{j,h},P_h$ for the corresponding finite graph triple and
product. At step $h$ apply the simultaneous lifting lemma to the three
prescribed graded maps
\begin{equation}\label{eq:threeK}
 \begin{array}{c|cc}
 \text{channel}&K_0(P_h)\to K_0(P_{h+1})&K_1(P_h)\to K_1(P_{h+1})\\\hline
 +&+1&0\\
 -&-1&0\\
 0&0&0.
 \end{array}
\end{equation}
Denote the resulting positive common and factor multiplicities by
$T_c,S_c,R_c$, for $c\in\{+,-,0\}$. With $q=q_h$, choose target sizes
by
\begin{equation}\label{eq:dimensions}
 \binom{k_{h+1}}{\ell_{h+1}}
       =(T_++T_-+qT_0)\binom{k_h}{\ell_h}.
\end{equation}
For the zero channel, use the source triple
$M_q(D_h),M_q(F_{0,h}),M_q(F_{1,h})$, whose graph weights are
$qk_h,q\ell_h$. Its multiplicities are $T_0,S_0,R_0$; hence its
contribution to the target dimensions is exactly $qT_0$.
The positive and negative channels use the unamplified source.

Place the three common representations in disjoint coordinate blocks
inside $D_{h+1}$. Their units are orthogonal constant projections
$p_+,p_-,p_0$ adding to one. For each factor, representation equivalence
of the common algebra gives a permutation unitary matching the entire
channel restriction, including the $M_q(D_h)$ restriction in channel
zero. Thus there are literal maps of triples with induced product maps
\[
 \gamma_+:P_h\to p_+P_{h+1}p_+,\quad
 \gamma_-:P_h\to p_-P_{h+1}p_-,\quad
 \widetilde\gamma_0:M_q(P_h)\to p_0P_{h+1}p_0.
\]
Here the full product of the amplified triple is $M_q(P_h)$.
This follows either by the common matrix units or by the same tensor
universal property below with coefficient $M_q$. All three source
products are simple, and each map is nonzero, hence injective.
The graded maps of $\widetilde\gamma_0$ under the usual corner Morita
identification are zero. They are not multiplied by $q$: the source
$K_0$ generator is a minimal projection of an amplified factor block.

Now tensor the triples by $E$. For any unital coefficient $E$ which is
nuclear, the full universal property gives
\begin{equation}\label{eq:tensor}
 (F_{0,h}\otimes E)*_{D_h\otimes E}(F_{1,h}\otimes E)
                  \cong Q_h:=P_h\otimes E.
\end{equation}
Indeed the common copy $1\otimes E$ commutes with both graph-factor
copies $F_{j,h}\otimes1$, hence with their full product. First obtain
the maximal tensor product and then use nuclearity of $E$. The reverse
map is defined on factor generators; both composites fix all generators.
The coefficient copy need not be internally commutative.

On the common and factor algebras, form the orthogonal sum of the
following maps, with the matching graph component in place of $\gamma$:
\begin{equation}\label{eq:Phi}
 \Phi_h(x)=
    (\gamma_+\otimes\id_E)(x)
   +(\gamma_-\otimes\eta\varepsilon)(x)
   +\bigl(\widetilde\gamma_0\circ(\id_{P_h}\otimes\sigma_h)\bigr)(x)
                                                        \otimes1_E.
\end{equation}
In the last term identify $P_h\otimes M_q$ with $M_q(P_h)$.
All channel units lie in the common algebra and are orthogonal, so the
sum is a unital *-homomorphism. The two constituent squares commute
literally, because each channel is a map of triples and its coefficient
map is the same on all three algebras. Thus \eqref{eq:Phi} is precisely
the induced map of full products, not a separately chosen map with
matching K-theory. The positive identity-coefficient channel is
injective (minimal tensor products preserve embeddings), so every
common, factor and product connecting map is injective.

\section{Exact KK calculation, including the non-UCT coefficient}
Each $P_h$ satisfies UCT: the finite-dimensional inclusions have
conditional expectations; the canonical two-sided cone is KK-equivalent
to $SP_h$ by \cite[Theorem 4.1]{FG}; its extension has finite-dimensional
quotient and suspended finite-dimensional ideal and therefore belongs
to the Rosenberg--Schochet bootstrap category. UCT transfers under
KK-equivalence and suspension by \cite[Theorem 1.17, Proposition 7.1]{RS}.
Together with $K_*(P_h)=(\Z,\Z)$, this gives a KK-equivalence
\[
                 \kappa_h:P_h\simeq_{KK}\C\oplus S\C
\]
respecting our chosen generators. One can construct inverse classes
using the UCT: its Ext term vanishes for free source K-groups, so the
prescribed inverse K-maps lift uniquely and their composites are identities.
The same observation shows that each even graph KK class is determined
by its two integer K-maps. There are no even off-diagonal entries between
$\C$ and $S\C$.

Exterior tensoring $\kappa_h$ with $1_E$ identifies $Q_h$ in KK with
\[
                            X=E\oplus SE.
\]
For the positive channel, the class is $\operatorname{diag}(1_E,0)$.
For the negative channel, it is
$\operatorname{diag}(-[\eta\varepsilon],0)$.
The graph class of $\widetilde\gamma_0$ is zero in KK, since its source
is Morita equivalent to the UCT algebra $P_h$ with free K-groups and
its two K-maps vanish. Its composition with any $\id\otimes\sigma_h$
and scalar inclusion is consequently KK-zero, even though it is an
injective actual homomorphism. Hence the total map in these coordinates is
\begin{equation}\label{eq:idempotent}
                   [\Phi_h]=f:=\operatorname{diag}(e,0),
                   \qquad e=1_E-[\varepsilon]\otimes[\eta].
\end{equation}
This argument uses no UCT or K\"unneth theorem for $E$ or $B$.

Let $L=\varinjlim(Q_h,\Phi_h)$. We justify fully why $L\simeq_{KK}F$.
In the KK category define $j:F\to X=E\oplus SE$ by $[\iota]$ in the
first coordinate and zero in the second, and $s:X\to F$ by $r$ in the
first coordinate. Then
\[
                  j\otimes s=1_F,\qquad s\otimes j=f.
\]
For each separable $Z$, the inverse system $\KK^i(X,Z)$ has bonding map
precomposition with the fixed idempotent $f$. It splits as a constant
identity system on $\operatorname{im}(f^*)$ and a zero-bonding system on
$\ker(f^*)$. Thus
\[
 \varprojlim\KK^i(X,Z)\cong\KK^i(F,Z),\qquad
 \varprojlim{}^1\KK^i(X,Z)=0.
\]
The second equality can also be seen directly from the product formula
for $\lim^1$: $1-$shift is surjective on the zero system and on the
constant identity system (solve recursively). No countability of these
KK groups or compactness of $F$ is being assumed.

All the $Q_h$ are separable nuclear, so the natural Milnor exact
sequence applies to $\KK^*(-,Z)$
\cite[Theorems 1.12 and 1.14(b)]{RS}. It gives
\begin{equation}\label{eq:naturaliso}
                        \KK^i(L,Z)\cong\KK^i(F,Z).
\end{equation}
Here is an explicit representing class for this isomorphism. Transport
$j$ to $j_h\in\KK(F,Q_h)$ using $\kappa_h\otimes1_E$.
Equation \eqref{eq:idempotent} implies
$j_h\otimes[\Phi_h]=j_{h+1}$. If $\lambda_h:Q_h\to L$ is the canonical
map, then $\beta=j_h\otimes[\lambda_h]\in\KK(F,L)$ is independent of $h$.
The isomorphism \eqref{eq:naturaliso} is precomposition with $\beta$.
Surjectivity with $Z=F$ gives $\alpha\in\KK(L,F)$ satisfying
$\beta\otimes\alpha=1_F$; injectivity with $Z=L$ gives
$\alpha\otimes\beta=1_L$. Hence $\alpha$ is an invertible class.
Moreover the restrictions of $\alpha$ to $Q_h$, expressed in the
$E\oplus SE$ coordinates, are exactly $s$. This follows either from
Milnor injectivity or from $s\otimes j=f$ and the compatibility equations.
Thus the retained K-class of an initial element is its reduced $F$
coordinate in the first $E$ summand. We will use this explicit statement
to prescribe the unit of a common corner.

\section{Simplicity, pure infiniteness, and stable finiteness}
Let $D_\infty,A_{0,\infty},A_{1,\infty}$ be the constituent limits.
They are separable nuclear and unital, and $D_\infty\to A_{j,\infty}$
is injective. The product limit $L$ is separable nuclear unital.
All connecting maps are injective, so we regard stages as subalgebras.

We first record two details about the evaluation mechanism. A channel
path in any iterated map has a support projection in the constant
finite-dimensional common algebra of the target stage. These supports
are orthogonal; compression by the selected support isolates that path.
This follows by induction from the orthogonal common supports at each
step: coefficient maps send constant projections to constant projections.
Second, a path of positive identity channels commutes with coefficient
evaluation by every $\sigma$. Thus, if a finite-dimensional representation
$\sigma$ detects a given nonzero positive stage element, following any
number of identity channels and then the zero/noise channel using that
same $\sigma$ still gives a nonzero element. Injectivity of the evaluated
graph or finite-dimensional maps proves this assertion.

Let $I$ be a nonzero ideal in any one of the four limits. Choose a
nonzero $b\in I_+$ and a positive stage approximation $a$, close enough
that for some $\epsilon>\|a-b\|$ the cut
$c=(a-\epsilon)_+$ is nonzero. Its image in the quotient by $I$ is zero,
so $c\in I$. The separating family supplies a $\sigma$ with
$(\id\otimes\sigma)(c)\ne0$. Select a later occurrence of that
representation in the tail sequence and follow identity channels until
then, and the zero/noise channel at that step. Compressing by the
constant path support gives a nonzero constant positive element $d\in I$.

For a constituent, at least one source matrix block evaluates nontrivially.
Every entry of $T_0,S_0,R_0$ is strictly positive. Consequently the
image $d$ is nonzero in every target constant matrix summand. A nonzero
positive matrix generates its matrix block as an ideal; hence $I$
contains every target constant block unit and therefore the common unit.
For the product, $d$ is a nonzero positive element of the simple constant
graph algebra $P_{h+1}$, so it generates that algebra's unit as an ideal.
This also forces $I=L$. This proves simplicity of all four limits.

For pure infiniteness of $L$, take any $0\ne b\in L_+$ and choose a
positive stage approximation $a$ and a nonzero cut
$c=(a-\epsilon)_+\precsim b$. The preceding evaluation procedure isolates
a nonzero constant positive $d$ at a later stage. It is an orthogonal
summand of the image of $c$, so $d\precsim c\precsim b$. The constant
graph algebra is simple purely infinite, so $\Her(d)$ there contains an
infinite projection $v$. This remains infinite in $L$, since embeddings
preserve the proper-isometry witness and its nonzero defect projection.
For completeness, from $v\precsim b$ choose $x$ so that $vx^*bxv$ is
invertible in $vLv$. Then
\[
         w=b^{1/2}xv(vx^*bxv)^{-1/2}
\]
has $w^*w=v$ and $ww^*\in\Her(b)$. Thus $\Her(b)$ contains an infinite
projection, as required. This proves pure infiniteness, without claiming
it for the coefficient tensor stages themselves.

Every unital RFD algebra has a tracial state, by composing a nonzero
unital finite-dimensional representation with normalized matrix trace.
Thus the unital constituent stages, finite sums of matrix algebras over
$E$, have nonempty compact trace spaces. For the inverse system of these
spaces, every finite compatibility requirement can be met by choosing a
trace at a sufficiently late stage and restricting backwards. Compactness
and the finite intersection property give compatible traces on each
limit, without any surjectivity assumption on trace restriction maps.
Simplicity makes the limit trace faithful. Its matrix traces are faithful,
so no matrix algebra over the constituent admits a proper isometry.
All three constituents are stably finite.

Finally, the literal diagrams imply the full universal identity
\begin{equation}\label{eq:fulllimit}
             L\cong A_{0,\infty}*_{D_\infty}A_{1,\infty}.
\end{equation}
Compatible factor maps give a map from the displayed full product to
$L$. Conversely, the two stage factors map compatibly into the displayed
full product, giving compatible maps from $Q_h$ and therefore from $L$.
The two composites fix every stage-factor generator and are identities.
All maps preserve the common unit.

\section{Prescribing the unit inside the common algebra}
The preceding construction identifies $L$ with $B$ up to KK-equivalence,
but an arbitrary unit must still be obtained by an actual common corner.
Let $u\in K_0(F)$ be the inverse image of $[1_B]$ under $\xi_*$.
By the definition of nonunital $K_0$, its image in $K_0(E)$ can be written
\[
                  \iota_*(u)=[p_0]-[q_0],\qquad p_0,q_0\in M_t(E)
\]
with equal scalar ranks under $\varepsilon$. Stabilize so that the
matrices have the same size; add the same nonzero scalar projection to
both if necessary to ensure the resulting common projection below is
nonzero. Choose the initial weights $k_{1,1},k_{1,2}$ large enough to
contain these two matrix projections. Choose positive $\ell_{1,1},\ell_{1,2}$
so that
\[
                   k_{1,1}+\ell_{1,1}=k_{1,2}+\ell_{1,2}.
\]
Then $[1_{P_1}]=0$ in $\coker B_1=\Z$, though the algebra is unital.
All later weights are chosen by \eqref{eq:dimensions} as before.

Put the actual projection $p$ in $D_1\otimes E$ by placing $p_0$ in the
first $k$-block, $q_0$ in the second $k$-block, and zero elsewhere.
The inclusion into $P_1\otimes E$ sends these two blocks to the positive
and negative $K_0(P_1)$ generators. Thus, in the first $E$ coordinate
of the KK identification $Q_1\simeq E\oplus SE$, its class is
$[p_0]-[q_0]=\iota_*(u)$, and its second coordinate is zero. This follows
also by external products with the two minimal graph projections; their
even KK classes are $+1$ and $-1$ in the scalar component and zero in
the suspension component. The calculation does not assume a K\"unneth
formula for $E$.

The KK equivalence $\alpha:L\to F$ constructed above consequently sends
$[p]$ to $u$. Since all maps are injective, $p$ is a nonzero projection
of the simple $D_\infty$, hence full there. It is also full in the other
three simple limit algebras.

\begin{lemma}[Common full corner]\label{lem:corner}
If $P=A_0*_D A_1$ and $p\in D$ is a full projection, then
\[
                    pPp\cong(pA_0p)*_{pDp}(pA_1p)
\]
as a full unital amalgam with unit $p$.
\end{lemma}
\begin{proof}
Fullness supplies finitely many $x_r\in Dp$ with
$\sum_rx_rx_r^*=1$: normalize a finite invertible positive sum from the
ideal generated by $p$. Set $R=(x_r^*x_s)\in M_N(pDp)$, a projection.
The maps $a\mapsto(x_r^*ax_s)$ identify $A_j$ with
$RM_N(pA_jp)R$ and agree on $D$. If $T=(pA_0p)*_{pDp}(pA_1p)$,
the full property gives $\Psi:P\to RM_N(T)R$. The column
$z=(x_r^*p)_r$ satisfies $z^*z=p$ and $zz^*=\Psi(p)$.
Thus $a\mapsto z^*\Psi(a)z$ maps $pPp$ to $T$ and fixes its two factor
corners. It is inverse to the natural map $T\to pPp$. To verify
generation of $pPp$ by these corners, insert
$1=\sum_rx_rx_r^*$ between successive letters of each compressed
alternating word; every resulting compressed factor lies in the
appropriate $pA_jp$. The dense word algebra proves both composites are
identities.
\end{proof}

Apply the lemma to \eqref{eq:fulllimit}. The three corners are separable
nuclear simple unital stably finite, and $pLp$ is a unital Kirchberg
algebra. Its full-corner inclusion into $L$ is a KK-equivalence. Compose
that class with $\alpha$ and $\xi$ to obtain an invertible element of
$\KK(pLp,B)$ taking its unit $[p]$ to $[1_B]$.
Phillips \cite[Corollary 4.2.2]{Phillips} gives a unital isomorphism
$pLp\cong B$. This precise corollary requires no UCT hypothesis.
The common-full-corner identity now proves Theorem \ref{thm:all}.

\section{Scope and verification status}
If correct, Theorem~\ref{thm:all} answers Suzuki's question for every
unital Kirchberg algebra, with nuclear constituents. It imposes no UCT
hypothesis and no restriction on the target's $K$-groups. The coefficient
algebra retains the target's full KK information; the argument does not
establish the UCT for arbitrary nuclear C*-algebras. It also does not
assert that the three constituents can always be chosen AF or AT.

The proof remains unverified. In particular, the RFD mapping-cone input,
the actual three-channel diagrams, the KK idempotent calculation, the
limit analysis, and the common-unit corner require independent scrutiny.
No claim of novelty or priority is established by this draft. Corrections
and comments are welcome.

\paragraph{Preparation and source note.}
This manuscript was developed with substantial AI assistance, including
drafting and automated review. Those reviews do not constitute independent
human mathematical verification. The cited statements of Rosenberg--Schochet
were consulted through text extracted from the original journal scan;
the complete original PDF was not available for direct inspection.
The bibliography identifies the source versions used, and theorem and page
numbers refer to those versions.

\begin{thebibliography}{9}
\bibitem{Suzuki} Y. Suzuki, \emph{Free groups amenably act on unital simple
AF-algebras}, \href{https://arxiv.org/abs/2610.01460v1}{arXiv:2610.01460v1},
1 October 2026. Corollary 1.4 and Remark 1.5.
\bibitem{KPR} A. Kumjian, D. Pask and I. Raeburn,
\emph{Cuntz--Krieger algebras of directed graphs}, February 1998 author
version, \href{https://documents.uow.edu.au/~dpask/index_files/papers/ckaodg.pdf}{author PDF}.
Theorem 1.2 and Corollary 3.11 (PDF p.~14). Published in Pacific J. Math.
\textbf{184} (1998), 161--174.
\bibitem{AG} P. Ara and K. R. Goodearl, \emph{C*-algebras of separated
graphs}, \href{https://arxiv.org/abs/1102.4296v2}{arXiv:1102.4296v2},
11 July 2011. Remark 3.10 records nuclearity of ordinary row-finite graph
algebras; separated-graph nuclearity is not asserted.
\bibitem{FG} P. Fima and E. Germain, \emph{The KK-theory of amalgamated
free products}, \href{https://arxiv.org/abs/1510.02418v3}{arXiv:1510.02418v3},
27 December 2016. Introduction p.~2; Theorem 4.1 p.~20; cone extension p.~25.
\bibitem{RS} J. Rosenberg and C. Schochet, \emph{The K\"unneth theorem
and the universal coefficient theorem for Kasparov's generalized K-functor},
Duke Math. J. \textbf{55} (1987), 431--474.
\href{https://doi.org/10.1215/S0012-7094-87-05524-4}{DOI}.
Theorems 1.12, 1.14(b), 1.17; Proposition 7.1. See the source note above.
\bibitem{Phillips} N. C. Phillips, \emph{A classification theorem for nuclear
purely infinite simple C*-algebras},
\href{https://arxiv.org/abs/funct-an/9506010v2}{arXiv:funct-an/9506010v2},
18 December 1997. Corollary 4.2.2, p.~42 (no UCT hypothesis).
\bibitem{Dadarlat} M. Dadarlat, \emph{Some remarks on the universal
coefficient theorem in KK-theory}, January 2002 author version,
\href{https://www.math.purdue.edu/~mdd/Papers.html}{author publication list}.
Lemma 2.4, pp.~2--3; construction in the proof of Theorem 1.2, p.~8.
Published in \emph{Operator Algebras and Mathematical Physics},
Theta (2003), pp.~65--74.
\end{thebibliography}
\end{document}
